\documentclass[11pt]{article}
\usepackage[margin=1in]{geometry}
\usepackage{lmodern}
\usepackage{booktabs}
\usepackage{hyperref}
\usepackage{enumitem}
\setlist{nosep,leftmargin=*}
\usepackage{xcolor}
\newcommand{\commentbox}[1]{\par\textit{#1}\par\medskip}
\newcommand{\changes}[1]{\textbf{Changes in the manuscript:} #1\par\medskip}

\title{Response to Reviewers\\[0.4em]
\large Training-Free Ranking from Pairwise Comparisons via Acyclic Graph Construction\\[0.2em]
\normalsize Journal of Supercomputing --- Major Revision}
\author{Soroush Vahidi}
\date{}

\begin{document}
\maketitle

\section*{Cover note to the Editor}

Dear Editor,

Thank you for arranging this careful review. We have prepared a major revision addressing all comments from Reviewers~1--4. The principal changes are:
\begin{itemize}
\item title revised to remove the word ``Scalable'';
\item corrected add-back implementation using exact reachability-aware cycle safety (replacing the submitted fixed-topological proxy);
\item secondary weighted min-cut structural exchange with proofs and regime-specific empirics;
\item rebuilt baseline comparisons from a single canonical leaderboard source (resolving the submitted Table~4/5 inconsistency);
\item direct classical runtime comparisons (showing that OURS is slower than lightweight classical methods);
\item comprehensive structural ablation, parameter sensitivity, timeout/coverage analysis, and family-aware statistics;
\item clarified ranking--MWFAS theory and DF03 guarantee boundaries;
\item manuscript-wide rewriting with a dedicated non-repetition pass.
\end{itemize}
No new scientific experiments remain open. The point-by-point responses below cite the revised manuscript sections, propositions, tables, and figures.

\vspace{0.8em}
\noindent Sincerely,\\
Soroush Vahidi

\section*{General reply to all reviewers}

Dear Editor and Reviewers,

Thank you for the careful and constructive reviews. We agree that the original presentation overstated novelty relative to prior local-ratio/add-back lineage, left the Demetrescu--Finocchi guarantee ambiguous under practical timeouts, and did not provide the ablation, classical runtime, and statistical evidence needed for a major revision. We have substantially revised the manuscript accordingly. Throughout, we distinguish \emph{prior} algorithmic lineage, an \emph{implementation-fidelity} correction, \emph{new theory}, a \emph{secondary} min-cut repair, and an \emph{expanded} empirical study.

%%%%%%%%%%%%%%%%%%%%%%%%%%%%%%%%%%%%%%%%%%%%%%%%%%%%%%%%%%%%%%%%%%%%%%
\section*{Reviewer 1}
%%%%%%%%%%%%%%%%%%%%%%%%%%%%%%%%%%%%%%%%%%%%%%%%%%%%%%%%%%%%%%%%%%%%%%

\paragraph{Comment 1 (Positioning / novelty).}
\commentbox{The contribution relative to prior feedback-arc-set / ranking work was not clearly delineated; the paper appeared to present the pipeline as more novel than supported.}

\noindent\textbf{Response.}
We agree that the original positioning was insufficient. We completely rewrote the Introduction, restructured Related Work (including an explicit relation to Vahidi \& Koutis, arXiv:2412.16181), and added a compact novelty-separation table. We no longer present local-ratio cycle breaking, exact cycle-safe reinsertion, weight-ordered add-back, or ternary refinement as newly invented. The title no longer contains ``Scalable.''

\changes{Sections~1 and~1.2; Table~1; title page.}

\paragraph{Comment 2 (Ablation / sensitivity).}
\commentbox{Provide ablations of Phase~1 only, Phase~1 with reinsertion, and the full pipeline; report upset metrics and runtime; study scale/density and sensitivity to zero tolerance, reinsertion passes, and refinement iterations.}

\noindent\textbf{Response.}
We added a structural ablation with human-readable stage labels:
A0~=~Phase~A only; A1~=~Phase~A + submitted topological-proxy add-back; A2~=~Phase~A + exact reachability add-back; A4~=~A2 + refinement; A5/A6~=~optional min-cut on top of A2/A4.
On non-Finance common completions, exact reachability add-back (A2) improves \texttt{upset\_simple} versus Phase~A only (A0) with W/T/L $76/0/1$ ($n=77$; Holm $p=2.4\times 10^{-12}$), and versus the submitted topological proxy (A1) with $32/0/1$ ($n=33$). Legacy multipass topo-proxy passes P2/P3 add little beyond P1; \texttt{zero\_tol} is stable; refinement saturates quickly (R1$\approx$R2). Non-Finance A4 runtimes are roughly $0.01$--$1.2$\,s (median $\approx 0.57$\,s) for $n\le 602$; Finance is the large-dense boundary (Section~3.5). Direct baseline quality/runtime tables retain the archived \textsc{OURS\_MFAS} leaderboard column for archive continuity; reachability gains are reported in the ablation (Table~7).

\changes{Sections~3.1, 3.5--3.8; Table~7; Figures~1--2.}

\paragraph{Comment 3 (Algorithm readability).}
\commentbox{The algorithm description needs complete, readable pseudocode and explicit parameters (ties, tolerances, timeouts).}

\noindent\textbf{Response.}
We inserted a unified Algorithm~1 covering deterministic cycle search, residual updates, zero tolerance, forced-progress safeguards, wall-clock checks, exact reachability reinsertion, optional min-cut, refinement, and identity fallback, together with a consolidated parameter table.

\changes{Sections~2.3 and~2.8; Algorithm~1; Table~2.}

\paragraph{Comment 4 (Theory / complexity / guarantees).}
\commentbox{Clarify the ranking--MWFAS relationship, complexity claims, and whether the practical method inherits the Demetrescu--Finocchi approximation guarantee; provide an error discussion for fallback ordering.}

\noindent\textbf{Response.}
Proposition~1 proves optimum-value equivalence $\mathrm{OPT}_{\mathrm{rank}}=\mathrm{OPT}_{\mathrm{MWFAS}}$ (many-to-many, not a bijection). For the audited Phase~A implementation, the idealized full-completion complexity is $O(mn+m^{2})$. Section~2.9 separates the idealized local-ratio guarantee from the practical finite-budget pipeline: we do \emph{not} claim the DF03 guarantee for prematurely terminated runs. Proposition~4 shows that identity-order fallback admits no nontrivial general multiplicative bound.

\changes{Sections~2.2, 2.9--2.11; Propositions~1 and~4; Remark on DF03.}

\paragraph{Comment 5 (Applications / future work).}
\commentbox{Provide more concrete future directions and application context.}

\noindent\textbf{Response.}
We added dedicated Limitations and Future Work sections covering denser-graph theory, SCC/parallel engineering, candidate selection for min-cut, and application settings (preference aggregation, multi-source sensors, industrial ranking) as directions, not evaluated domains.

\changes{Sections~4--6.}

%%%%%%%%%%%%%%%%%%%%%%%%%%%%%%%%%%%%%%%%%%%%%%%%%%%%%%%%%%%%%%%%%%%%%%
\section*{Reviewer 2}
%%%%%%%%%%%%%%%%%%%%%%%%%%%%%%%%%%%%%%%%%%%%%%%%%%%%%%%%%%%%%%%%%%%%%%

\paragraph{Comment 1 (Novelty).}
\commentbox{Clarify what is scientifically new relative to prior MWFAS / ranking work.}

\noindent\textbf{Response.}
As in our reply to Reviewer~1 (Comment~1), we now separate prior lineage, fidelity correction, new theory, secondary min-cut, and expanded empirics (Table~1; Sections~1--2).

\paragraph{Comment 2 (Add-back contribution).}
\commentbox{The submitted add-back / INS variants do not appear to contribute clearly; please diagnose and strengthen this stage.}

\noindent\textbf{Response.}
We agree that the submitted fixed-topological proxy was too conservative and that multipass INS2/INS3 were not major quality innovations. Exact reachability restores the intended cycle-safety test and materially improves rankings in the ablation (A0$\to$A2: $76/0/1$; A1$\to$A2: $32/0/1$). INS multipass is retained only historically. As a secondary mechanism we added a weighted min-cut exchange with monotone weighted-FAS improvement (Proposition~3; Section~3.9). Direct baseline Tables~4--6 continue to report the archived topo-proxy \textsc{OURS\_MFAS} column for continuity; the revised canonical pipeline is the reachability-aware method of Section~2.

\changes{Sections~2.5--2.6, 3.7--3.9; Algorithms~1--2; Table~7.}

\paragraph{Comment 3 (GNN protocol / timing fairness).}
\commentbox{Clarify what GNN runtimes include and whether hardware is comparable.}

\noindent\textbf{Response.}
Archived GNN timings measure end-to-end training+evaluation, not inference-only latency. The exact archived GPU model ledger is unavailable; the manuscript therefore qualifies hardware comparability and does not claim matched stage-by-stage inference fairness.

\changes{Section~3.1; runtime discussion in Section~3.4.}

\paragraph{Comment 4 (Timeout bias / missingness).}
\commentbox{Timeouts and missing values may bias aggregates; please report coverage honestly.}

\noindent\textbf{Response.}
Pairwise common-completion comparisons are now primary. Coverage is reported explicitly (e.g., OURS $77/79$; classical typically $78/79$; DIGRAC/ib $61/79$). Finance is separated rather than averaged into quality macros. We do not impute NaNs as scores.

\changes{Sections~3.1, 3.4, 3.5.}

\paragraph{Comment 5 (Oracle best-in-suite).}
\commentbox{Best-in-suite / oracle envelopes should not be treated as deployable competitors.}

\noindent\textbf{Response.}
We agree. Direct fixed-method comparisons are primary; oracle envelopes are de-emphasized and labeled as non-deployable post-hoc context only.

\changes{Section~3.1; captions of Tables~4--5.}

\paragraph{Comment 6 (Statistics / dataset dependence).}
\commentbox{Provide formal paired statistics and address Basketball-dominated dependence.}

\noindent\textbf{Response.}
We report W/T/L, median paired deltas, and---where exported---Wilcoxon tests with Holm adjustment and bootstrap CIs. Family-aware equal-family macros, leave-one-family-out, and Basketball-collapsed analyses are included. SpringRank comparisons are acknowledged as LOFO-fragile; BTL's \texttt{upset\_ratio} advantage persists across families.

\changes{Sections~3.2--3.3; Tables~4--5,~7.}

\paragraph{Comment 7 (Deterministic repetitions).}
\commentbox{Ten deterministic repetitions should not be interpreted as ranking robustness.}

\noindent\textbf{Response.}
We agree. For deterministic non-GNN methods, repeats characterize runtime variability and timeout repeatability, not independent ranking observations.

\changes{Sections~3.1 and~4.}

\paragraph{Comment 8 (Scalability).}
\commentbox{The word ``scalable'' is not supported without qualification given dense/large cases.}

\noindent\textbf{Response.}
We removed ``Scalable'' from the title and qualify empirical scope to the evaluated sparse and moderately dense graphs in our suite, with Finance as an explicit boundary.

\changes{Title; Sections~3.5--3.6, 4.}

\paragraph{Comment 9 (Presentation / repetition).}
\commentbox{The manuscript is repetitive and hard to navigate.}

\noindent\textbf{Response.}
We consolidated the experimental protocol, separated Limitations from Conclusion, and performed a manuscript-wide repetition audit (novelty, DF03, INS, Finance, runtime accounting, contributions).

\changes{Throughout; Sections~3.1, 4--5.}

%%%%%%%%%%%%%%%%%%%%%%%%%%%%%%%%%%%%%%%%%%%%%%%%%%%%%%%%%%%%%%%%%%%%%%
\section*{Reviewer 3}
%%%%%%%%%%%%%%%%%%%%%%%%%%%%%%%%%%%%%%%%%%%%%%%%%%%%%%%%%%%%%%%%%%%%%%

\paragraph{Comment 1 (Ineffective add-back / INS1--3).}
\commentbox{OURS-MFAS and INS variants appear nearly identical; add-back does not seem to help.}

\noindent\textbf{Response.}
We agree with this diagnosis of the submitted manuscript: archived \textsc{OURS-MFAS} and \textsc{INS3} were effectively duplicate topo-proxy configurations, and the fixed-topological rule was too conservative. The revised Phase~B uses exact reachability-aware reinsertion. Empirically, A2 improves \texttt{upset\_simple} over Phase~A on $76/77$ datasets and over the legacy proxy on $32/33$. Multipass INS2/INS3 are de-emphasized; reachability (and optional min-cut) replace them as the corrected/structural path.

\changes{Sections~2.5, 3.1, 3.7--3.8; Table~7.}

\paragraph{Comment 2 (Table~4 / Table~5 inconsistency).}
\commentbox{Submitted baseline tables were inconsistent (including classical values that did not match).}

\noindent\textbf{Response.}
We agree and take responsibility for this error. Stale aggregation/pivot logic produced inconsistent submitted summaries. The revision rebuilds all baseline tables from one canonical per-method leaderboard export under pairwise common-completion rules. For example, the canonical SpringRank median \texttt{upset\_simple} in the verified full-suite export is $0.802724$. Submitted unpaired Table~4/5 aggregates are not reused.

\changes{Section~3.1 (canonical source statement); Tables~4--5.}

\paragraph{Comment 3 (Runtime versus classical methods).}
\commentbox{Compare runtime directly against classical baselines, not only against GNNs.}

\noindent\textbf{Response.}
We now report paired runtime W/T/L against classical methods. OURS is slower than lightweight classical estimators (median ratios from about $2.6\times$ vs SyncRank to $>500\times$ vs PageRank). The runtime advantage is only relative to archived trained GNN end-to-end runs ($\sim 8\times$ faster on $60/60$ common completions vs DIGRAC/ib), with the protocol qualification noted above.

\changes{Section~3.4; Table~6.}

\paragraph{Comment 4 (Accuracy--runtime tradeoff / overclaiming).}
\commentbox{Claims of accuracy and runtime advantage need to be metric-specific and honest.}

\noindent\textbf{Response.}
We revised claims to be metric-specific: under the archived \textsc{OURS\_MFAS} column used for Tables~4--5, strong simple/naive upset performance against several baselines; competitive (not uniformly superior) with SpringRank/DavidScore/SVD\_NRS under family-aware analysis; BTL remains stronger on \texttt{upset\_ratio}. Ablations separately show that restoring exact reachability improves \texttt{upset\_simple} relative to Phase~A and the topo proxy. We do not claim universal dominance. Together with Comments~1--3, the three concrete requests---reconcile tables, fix/strengthen insertion, and add classical runtime comparisons---are implemented.

\changes{Abstract; Sections~3.2--3.4, 5; Tables~4--6.}

%%%%%%%%%%%%%%%%%%%%%%%%%%%%%%%%%%%%%%%%%%%%%%%%%%%%%%%%%%%%%%%%%%%%%%
\section*{Reviewer 4}
%%%%%%%%%%%%%%%%%%%%%%%%%%%%%%%%%%%%%%%%%%%%%%%%%%%%%%%%%%%%%%%%%%%%%%

\paragraph{Comment 1 (Scientific novelty).}
\commentbox{What exactly is scientifically new?}

\noindent\textbf{Response.}
The Introduction and Related Work now answer this directly via prior/new separation and Table~1: formal ranking--MWFAS clarification and guarantee boundaries; fidelity restoration of exact add-back; secondary min-cut exchange; expanded empirical evaluation. Prior local-ratio/add-back/refinement components are credited, not reinvented.

\changes{Sections~1--2; Table~1.}

\paragraph{Comment 2 (DF03 approximation guarantee).}
\commentbox{Does the practical implementation inherit the Demetrescu--Finocchi guarantee?}

\noindent\textbf{Response.}
We clarify that the published guarantee applies to the idealized local-ratio procedure under its completion assumptions. The practical pipeline adds wall-clock termination, floating-point tolerances, forced-progress safeguards, and fallback output; therefore prematurely terminated practical runs do \emph{not} automatically inherit the DF03 guarantee (Remark in Section~2.9). Proposition~4 addresses fallback error.

\changes{Sections~2.9--2.10, 4.}

\paragraph{Comment 3 (Multipass INS).}
\commentbox{INS1/INS2/INS3 appear oversold.}

\noindent\textbf{Response.}
Sensitivity analysis shows P1 restores many edges while P2/P3 add little ranking quality. Multipass INS variants are no longer framed as a core contribution.

\changes{Sections~2 (opening), 3.1, 3.8.}

\paragraph{Comment 4 (GNN protocol).}
\commentbox{Clarify GNN timing/accounting.}

\noindent\textbf{Response.}
As in our reply to Reviewer~2 (Comment~3): end-to-end trained-GNN protocol, not inference-only; hardware ledger unavailable; claims qualified in Section~3.1.

\paragraph{Comment 5 (``Scalable'').}
\commentbox{The title/text overstate scalability.}

\noindent\textbf{Response.}
We removed ``Scalable'' from the title and describe empirical behavior on the evaluated sparse and moderately dense graphs in our suite, with Finance as an explicit large-dense boundary.

\changes{Title; Sections~3.5--3.6, 4.}

\paragraph{Comment 6 (Writing / repetition).}
\commentbox{The writing is repetitive and sentences are too long.}

\noindent\textbf{Response.}
We rewrote the Introduction, consolidated protocol and limitations, removed duplicated novelty/DF03/INS/Finance explanations, and shortened long passages while preserving technical precision.

\changes{Throughout; Sections~1, 3.1, 4--5.}

\paragraph{Comment 7 (MWFAS backbone / cycle-selection ablations).}
\commentbox{Ablate the MWFAS backbone and related algorithmic choices.}

\noindent\textbf{Response.}
The A0--A6 structural ablation addresses backbone stages (Phase~A only; topo-proxy add-back; exact reachability; refinement; optional min-cut). Cycle-selection sensitivity (C0/C1) shows a small but detectable effect on the A4 path ($|\mathrm{med}\,\Delta|\approx 10^{-3}$); we report it without overstatement.

\changes{Sections~3.7--3.8; Table~7; Figure~2.}

\section*{Closing}

We thank the reviewers again for comments that substantially improved the clarity, correctness, and evidence standards of the manuscript. We hope the revised version addresses all concerns.

\end{document}
